\documentclass[11pt,a4paper]{article}
\usepackage[top=2.3cm,bottom=2.3cm,left=2.2cm,right=2.2cm]{geometry}

\usepackage{amsmath,amssymb}
\usepackage{fontspec}
\setmainfont{Liberation Serif}
\setsansfont{Liberation Sans}
\setmonofont{DejaVu Sans Mono}[Scale=0.84]

\usepackage{xcolor}
\definecolor{brandblue}{RGB}{20, 55, 105}
\definecolor{accentblue}{RGB}{35, 95, 165}
\definecolor{codebg}{RGB}{248, 249, 251}
\definecolor{codeframe}{RGB}{215, 222, 230}
\definecolor{javakeyword}{RGB}{0, 45, 170}
\definecolor{javastring}{RGB}{5, 120, 20}
\definecolor{javacomment}{RGB}{120, 120, 120}
\definecolor{javanumber}{RGB}{20, 75, 220}

\definecolor{noteborder}{RGB}{30, 110, 65}
\definecolor{notebg}{RGB}{242, 249, 245}
\definecolor{warnborder}{RGB}{185, 75, 10}
\definecolor{warnbg}{RGB}{254, 248, 240}
\definecolor{tipborder}{RGB}{30, 95, 160}
\definecolor{tipbg}{RGB}{242, 247, 254}

\usepackage{fancyhdr}
\usepackage{lastpage}
\pagestyle{fancy}
\fancyhf{}
\lhead{\parbox[t]{0.58\textwidth}{\raggedright\small\sffamily\color{brandblue}\textbf{T05-T06 --- Sintassi Java, Strutture di Controllo e Classi Base}}}
\rhead{\small\sffamily\color{gray}Programmazione II -- UniVR (2025/2026)}
\cfoot{\small\sffamily Pagina \thepage\ di \pageref{LastPage}}
\renewcommand{\headrulewidth}{0.5pt}
\renewcommand{\headrule}{\hbox to\headwidth{\color{brandblue!70!black}\leaders\hrule height \headrulewidth\hfill}}

\usepackage{titlesec}
\titleformat{\section}{\Large\bfseries\sffamily\color{brandblue}}{\thesection}{0.8em}{}[\titlerule]
\titleformat{\subsection}{\large\bfseries\sffamily\color{accentblue}}{\thesubsection}{0.8em}{}
\titleformat{\subsubsection}{\normalsize\bfseries\sffamily\color{brandblue!85!black}}{\thesubsubsection}{0.8em}{}

\usepackage{needspace}
\usepackage{fancyvrb}
\DefineVerbatimEnvironment{asciibox}{Verbatim}{
  fontsize=\fontsize{7.5pt}{9.2pt}\selectfont,
  frame=single,
  rulecolor=\color{codeframe},
  fillcolor=\color{codebg},
  framesep=2.5mm
}
\usepackage{listings}
\lstset{
  backgroundcolor=\color{codebg},
  frame=single,
  rulecolor=\color{codeframe},
  basicstyle=\ttfamily\footnotesize,
  keywordstyle=\color{javakeyword}\bfseries,
  stringstyle=\color{javastring},
  commentstyle=\color{javacomment}\itshape,
  numberstyle=\tiny\color{gray},
  numbers=left,
  numbersep=7pt,
  stepnumber=1,
  breaklines=true,
  breakatwhitespace=true,
  showstringspaces=false,
  tabsize=4,
  xleftmargin=12pt,
  framexleftmargin=12pt
}

\newcommand{\passthrough}[1]{#1}
\providecommand{\tightlist}{\setlength{\itemsep}{0pt}\setlength{\parskip}{0pt}}

\usepackage{tcolorbox}
\tcbuselibrary{skins,breakable}
\newtcolorbox{studynote}[1][Nota]{
  enhanced,
  breakable,
  colback=notebg,
  colframe=noteborder,
  fonttitle=\bfseries\sffamily\small,
  coltitle=white,
  title={#1},
  arc=2mm,
  boxrule=0.8pt,
  left=9pt, right=9pt, top=7pt, bottom=7pt
}

\newtcolorbox{studywarning}[1][Attenzione]{
  enhanced,
  breakable,
  colback=warnbg,
  colframe=warnborder,
  fonttitle=\bfseries\sffamily\small,
  coltitle=white,
  title={#1},
  arc=2mm,
  boxrule=0.8pt,
  left=9pt, right=9pt, top=7pt, bottom=7pt
}

\newtcolorbox{studytip}[1][Suggerimento]{
  enhanced,
  breakable,
  colback=tipbg,
  colframe=tipborder,
  fonttitle=\bfseries\sffamily\small,
  coltitle=white,
  title={#1},
  arc=2mm,
  boxrule=0.8pt,
  left=9pt, right=9pt, top=7pt, bottom=7pt
}

\usepackage{booktabs}
\usepackage{longtable}
\usepackage{array}
\usepackage{calc}

\usepackage{hyperref}
\hypersetup{
  colorlinks=true,
  linkcolor=accentblue,
  urlcolor=accentblue,
  citecolor=brandblue,
  pdfborder={0 0 0}
}

\title{\Huge\bfseries\sffamily\color{brandblue} T05–T06 --- Sintassi Java, Strutture di Controllo e Classi Base \\[0.3em] \large\sffamily Costrutti Condizionali e Iterativi, Definizione di Classe e Nascita del Progetto SimpleDate}
\author{\textbf{Docente:} Prof. Michele Pasqua \quad---\quad \textbf{Appunti Revisione:} Wally}
\date{A.A. 2025/2026 -- Universit\`a degli Studi di Verona}

\begin{document}
\maketitle
\thispagestyle{fancy}
\vspace{-1.2em}
\noindent\textcolor{brandblue}{\rule{\textwidth}{0.8pt}}
\vspace{1.2em}

In questo blocco analizziamo in dettaglio la sintassi fondamentale di
Java
(\href{file:///home/wally/Documenti/GitHub/Programmazione-II/25-26/Teoria-e-Esercitazioni/Slides-20260902/T05\%20-\%20Java\%20Basic\%20Syntax.pdf}{T05
- Java Basic Syntax.pdf}), la struttura delle classi e degli oggetti
(\href{file:///home/wally/Documenti/GitHub/Programmazione-II/25-26/Teoria-e-Esercitazioni/Slides-20260902/T06\%20-\%20Java\%20Classes\%20and\%20Objects.pdf}{T06
- Java Classes and Objects.pdf}) e il codice sorgente di debutto del
corso: il progetto
\href{file:///home/wally/Documenti/GitHub/Programmazione-II/25-26/Teoria-e-Esercitazioni/Codice-20260902/Lezioni05-06/SimpleDate}{\passthrough{\lstinline!SimpleDate!}}
e il file di laboratorio
\href{file:///home/wally/Documenti/GitHub/Programmazione-II/25-26/Teoria-e-Esercitazioni/Codice-20260902/Lezioni05-06/StringPlayground.java}{\passthrough{\lstinline!StringPlayground.java!}}.

\begin{center}\rule{0.5\linewidth}{0.5pt}\end{center}

\subsubsection{1. Teoria: Concetti Teorici dalle
Slide}\label{teoria-concetti-teorici-dalle-slide}

\paragraph{A. Sintassi di Base e Costrutti Semplici (Slide T05,
1--19)}\label{a.-sintassi-di-base-e-costrutti-semplici-slide-t05-119}

\begin{itemize}
\tightlist
\item
  \textbf{Variabili e Inizializzazione Automatica}:

  \begin{itemize}
  \tightlist
  \item
    Una variabile associa un nome e un tipo a una locazione di memoria.
  \item
    \textbf{Regola cruciale}: Solo i campi di classe/istanza
    (\emph{fields}) vengono inizializzati automaticamente al default
    (\passthrough{\lstinline!0!}, \passthrough{\lstinline!0.0!},
    \passthrough{\lstinline!false!},
    \passthrough{\lstinline!'\\u0000'!},
    \passthrough{\lstinline!null!}). Le variabili locali nello Stack
    richiedono l'assegnamento esplicito prima dell'uso.
  \end{itemize}
\item
  \textbf{Costanti (\passthrough{\lstinline!static final!})}:

  \begin{itemize}
  \tightlist
  \item
    \passthrough{\lstinline!final static int I = 5;!}: la variabile è a
    sola lettura; deve essere inizializzata contestualmente o nel
    costruttore e non può più essere riassegnata.
  \end{itemize}
\item
  \textbf{Identificatori e Convenzioni di Naming}:

  \begin{itemize}
  \tightlist
  \item
    Regola sintattica: sequenza arbitraria di caratteri
    \passthrough{\lstinline![a-z]!}, \passthrough{\lstinline![A-Z]!},
    \passthrough{\lstinline![0-9]!}, \passthrough{\lstinline!\_!} e
    \passthrough{\lstinline!$!}, che non inizi con una cifra e non
    coincida con una delle \textbf{50 parole chiave riservate}
    (\passthrough{\lstinline!class!}, \passthrough{\lstinline!static!},
    \passthrough{\lstinline!new!}, \passthrough{\lstinline!this!},
    \passthrough{\lstinline!super!}, \passthrough{\lstinline!switch!},
    \passthrough{\lstinline!goto!}, \passthrough{\lstinline!const!},
    ecc.).
  \item
    Convenzioni ufficiali Java:

    \begin{itemize}
    \tightlist
    \item
      \textbf{Classi}: \passthrough{\lstinline!UpperCamelCase!} (es.
      \passthrough{\lstinline!SimpleDate!},
      \passthrough{\lstinline!MainDate!}).
    \item
      \textbf{Metodi, variabili locali, campi e parametri}:
      \passthrough{\lstinline!lowerCamelCase!} (es.
      \passthrough{\lstinline!daysPerMonth!},
      \passthrough{\lstinline!myCar!}).
    \item
      \textbf{Costanti}: \passthrough{\lstinline!UPPER\_SNAKE\_CASE!}
      (es. \passthrough{\lstinline!FORMAT!},
      \passthrough{\lstinline!PI\_CONST!}).
    \end{itemize}
  \end{itemize}
\item
  \textbf{Caratteri di Escape Speciali}:

  \begin{itemize}
  \tightlist
  \item
    \passthrough{\lstinline!\\n!} (LF linefeed,
    \passthrough{\lstinline!\\u000a!}), \passthrough{\lstinline!\\r!}
    (CR carriage return, \passthrough{\lstinline!\\u000d!}),
    \passthrough{\lstinline!\\t!} (tab orizzontale,
    \passthrough{\lstinline!\\u0009!}), \passthrough{\lstinline!\\b!}
    (backspace, \passthrough{\lstinline!\\u0008!}),
    \passthrough{\lstinline!\\f!} (form feed,
    \passthrough{\lstinline!\\u000c!}), \passthrough{\lstinline!\\"!}
    (doppio apice), \passthrough{\lstinline!\\'!} (singolo apice),
    \passthrough{\lstinline!\\\\!} (backslash).
  \end{itemize}
\item
  \textbf{Commenti e Documentazione}:

  \begin{itemize}
  \tightlist
  \item
    \passthrough{\lstinline!// ...!} (singola riga).
  \item
    \passthrough{\lstinline!/* ... */!} (multi-riga).
  \item
    \passthrough{\lstinline!/** ... */!} (\textbf{Javadoc}): elaborato
    dal tool \passthrough{\lstinline!javadoc!} per generare
    automaticamente la documentazione HTML delle API.
  \end{itemize}
\item
  \textbf{Espressioni, Valutazione Sinistra-Destra e Side Effects}:

  \begin{itemize}
  \tightlist
  \item
    In Java le sotto-espressioni sono valutate \textbf{rigorosamente da
    sinistra a destra}:

    \begin{itemize}
    \tightlist
    \item
      \passthrough{\lstinline!int i = 1; int j = (i = 2) * i;!}
      \(\implies\) \passthrough{\lstinline!i!} diventa
      \passthrough{\lstinline!2!}, poi moltiplicato per
      \passthrough{\lstinline!i!} (\passthrough{\lstinline!2!})
      \(\implies\) \passthrough{\lstinline!j = 4!}.
    \item
      \passthrough{\lstinline!int t = 3; eval(t + (t = 2))!}
      \(\implies\) prima si valuta \passthrough{\lstinline!t!}
      (\passthrough{\lstinline!3!}), poi
      \passthrough{\lstinline!(t = 2)!} (\passthrough{\lstinline!2!})
      \(\implies 3 + 2 = 5\).
    \end{itemize}
  \end{itemize}
\item
  \textbf{Operatori Aritmetici Unari: Pre-incremento vs
  Post-incremento}:

  \begin{itemize}
  \item
    \passthrough{\lstinline!x++!} (post-incremento): restituisce il
    valore attuale di \passthrough{\lstinline!x!} e
    \emph{successivamente} incrementa \passthrough{\lstinline!x!} di 1.
  \item
    \passthrough{\lstinline!++x!} (pre-incremento): incrementa
    \emph{subito} \passthrough{\lstinline!x!} di 1 e restituisce il
    nuovo valore.
  \item
    \textbf{L'esercizio tranello delle slide}:

\begin{lstlisting}[language=Java]
int t = 3;
t = -t++ - ++t; // Quanto vale t?
\end{lstlisting}

    \begin{itemize}
    \tightlist
    \item
      Valutazione da sinistra a destra:

      \begin{enumerate}
      \def\labelenumi{\arabic{enumi}.}
      \tightlist
      \item
        \passthrough{\lstinline!-t++!}: valuta \(-3\), e incrementa
        \(t\) a \(4\).
      \item
        \passthrough{\lstinline!++t!}: incrementa subito \(t\) da \(4\)
        a \(5\), e valuta \(5\).
      \item
        Operazione: \((-3) - 5 = -8\). Assegnato a \(t\), il valore
        finale è \textbf{\passthrough{\lstinline!-8!}}.
      \end{enumerate}
    \end{itemize}
  \end{itemize}
\item
  \textbf{Operatore Ternario}:
  \passthrough{\lstinline!condizione ? expr1 : expr2!}.
\item
  \textbf{Operatori Logici e Corto-Circuito}:

  \begin{itemize}
  \tightlist
  \item
    \passthrough{\lstinline!\&!}, \passthrough{\lstinline!|!}: valutano
    sempre entrambi i rami.
  \item
    \passthrough{\lstinline!\&\&!}, \passthrough{\lstinline!||!}:
    valutazione a corto circuito (\emph{short-circuit}). Se il primo
    operando determina già il risultato, il secondo non viene nemmeno
    eseguito (essenziale per prevenire crash su puntatori nulli).
  \end{itemize}
\item
  \textbf{Assegnamenti Composti}: \passthrough{\lstinline!x += expr!}
  equivale a \passthrough{\lstinline!x = x + (expr)!}. Se \(x=4\),
  l'espressione \(x += x - 7\) assegna \(x = 4 + (4 - 7) = 1\).
\end{itemize}

\begin{center}\rule{0.5\linewidth}{0.5pt}\end{center}

\paragraph{B. Costrutti di Controllo di Flusso (Slide T05,
20--29)}\label{b.-costrutti-di-controllo-di-flusso-slide-t05-2029}

\begin{itemize}
\tightlist
\item
  \textbf{Blocchi \passthrough{\lstinline!\{ ... \}!}}: delimitano lo
  scope di vita delle variabili locali.
\item
  \textbf{Condizionali}:
  \passthrough{\lstinline!if (Bexpr) ... else ...!}.
\item
  \textbf{Costrutto \passthrough{\lstinline!switch-case!} e il
  Fall-Through}:

  \begin{itemize}
  \tightlist
  \item
    Seleziona il ramo da eseguire in base al valore di un'espressione
    (\passthrough{\lstinline!byte!}, \passthrough{\lstinline!short!},
    \passthrough{\lstinline!char!}, \passthrough{\lstinline!int!},
    \passthrough{\lstinline!String!}, \passthrough{\lstinline!enum!}).
  \item
    Se non si inserisce l'istruzione \passthrough{\lstinline!break!},
    l'esecuzione \textbf{continua nei rami successivi}
    (\emph{fall-through}). Questo meccanismo, se usato volontariamente,
    permette di raggruppare casi che condividono la stessa logica (come
    vedremo in \passthrough{\lstinline!daysPerMonth!}).
  \end{itemize}
\item
  \textbf{Cicli}:

  \begin{itemize}
  \item
    \passthrough{\lstinline!while (cond) \{ ... \}!}: test preliminare
    (può eseguire 0 volte).
  \item
    \passthrough{\lstinline!do \{ ... \} while (cond);!}: test
    posticipato (esegue almeno 1 volta).
  \item
    \passthrough{\lstinline!for (init; cond; step) \{ ... \}!}: ciclo
    controllato. Supporta inizializzazioni e passi multipli separati da
    virgola:

\begin{lstlisting}[language=Java]
for (int i = 0, j = 9; i <= 9 && j >= 0; i++, j--) { ... }
\end{lstlisting}
  \end{itemize}
\item
  \textbf{Salto di Controllo}:

  \begin{itemize}
  \tightlist
  \item
    \passthrough{\lstinline!break!}: interrompe il ciclo o lo switch ed
    esce dal blocco.
  \item
    \passthrough{\lstinline!continue!}: salta immediatamente alla
    successiva iterazione del ciclo.
  \item
    \passthrough{\lstinline!return!}: esce dal metodo corrente
    restituendo eventualmente un valore.
  \end{itemize}
\end{itemize}

\begin{center}\rule{0.5\linewidth}{0.5pt}\end{center}

\paragraph{\texorpdfstring{C. Classi, Oggetti, \texttt{this} e
\texttt{toString} (Slide T06,
1--28)}{C. Classi, Oggetti, this e toString (Slide T06, 1--28)}}\label{c.-classi-oggetti-this-e-tostring-slide-t06-128}

\begin{itemize}
\tightlist
\item
  \textbf{Perché modellare una data? (La vignetta a lezione, Slide 7)}:

  \begin{itemize}
  \tightlist
  \item
    Lo slide show mostra il fumetto geek: \emph{«Imagine an object much
    like a thing in the real world, like\ldots{} A DATE?»}. È da qui che
    nasce la classe \passthrough{\lstinline!Date!} come filo conduttore
    del corso.
  \end{itemize}
\item
  \textbf{Anatomia di una Classe}:

  \begin{itemize}
  \tightlist
  \item
    \textbf{Campi d'istanza (\emph{Fields})}: variabili che memorizzano
    lo stato di ciascun oggetto.
  \item
    \textbf{Metodi d'istanza (\emph{Methods})}: funzioni che definiscono
    il comportamento e i messaggi accettati dall'oggetto.
  \item
    \textbf{Firma del metodo}:
    \passthrough{\lstinline!tipoRitorno nomeMetodo(tipoParametro nomeParametro, ...)!}.
    Se non restituisce nulla si usa \passthrough{\lstinline!void!}.
  \item
    \textbf{Overloading}: possibilità di definire più metodi con lo
    stesso nome purché differiscano per numero, tipo o ordine dei
    parametri.
  \end{itemize}
\item
  \textbf{Entrypoint \passthrough{\lstinline!main!} e Metodi Statici
  (\passthrough{\lstinline!static!})}:

  \begin{itemize}
  \tightlist
  \item
    \passthrough{\lstinline!public static void main(String[] args)!}: il
    launcher della JVM non passa il nome del programma in
    \passthrough{\lstinline!args[0]!} (a differenza del C).
    \passthrough{\lstinline!args[0]!} è già il primo parametro utente.
  \item
    \textbf{Metodi Statici}: appartengono alla classe e non alle singole
    istanze. Si invocano con
    \passthrough{\lstinline!NomeClasse.metodo()!}. Non possono accedere
    allo stato d'istanza né usare \passthrough{\lstinline!this!}.
  \end{itemize}
\item
  \textbf{Creazione dell'Oggetto (\passthrough{\lstinline!new!}) e
  Costruttore}:

  \begin{itemize}
  \tightlist
  \item
    L'istruzione \passthrough{\lstinline!new Date(...)!}:

    \begin{enumerate}
    \def\labelenumi{\arabic{enumi}.}
    \tightlist
    \item
      Alloca memoria nello Heap per l'oggetto.
    \item
      Inizializza i campi ai valori di default.
    \item
      Invoca il \textbf{costruttore} (metodo speciale con lo stesso nome
      della classe e senza tipo di ritorno) per valorizzare lo stato.
    \item
      Restituisce il riferimento all'oggetto appena creato.
    \end{enumerate}
  \end{itemize}
\item
  \textbf{Aliasing e Garbage Collection}:

  \begin{itemize}
  \tightlist
  \item
    Se scriviamo \passthrough{\lstinline!Date d2 = d1;!}, le due
    variabili puntano allo stesso identico oggetto nello Heap
    (\emph{aliasing}).
  \item
    Quando una reference viene azzerata
    (\passthrough{\lstinline!d1 = null; d2 = null;!}), l'oggetto diventa
    orfano e il Garbage Collector può bonificarlo.
  \end{itemize}
\item
  \textbf{La Parola Chiave \passthrough{\lstinline!this!}}:

  \begin{itemize}
  \item
    È il riferimento implicito all'oggetto corrente su cui è stato
    invocato il metodo.
  \item
    È \textbf{obbligatoria} quando i parametri del costruttore/metodo
    hanno lo stesso identico nome dei campi d'istanza (\emph{shadowing
    dei parametri}):

\begin{lstlisting}[language=Java]
public Date(int day, int month, int year) {
    this.day = day;     // this.day è il campo dell'oggetto nello Heap
    this.month = month; // day è il parametro locale nello Stack
    this.year = year;
}
\end{lstlisting}
  \end{itemize}
\item
  \textbf{Rappresentazione a Stringa e
  \passthrough{\lstinline!toString()!}}:

  \begin{itemize}
  \tightlist
  \item
    Tutte le classi Java ereditano dalla superclasse universale
    \passthrough{\lstinline!java.lang.Object!} il metodo
    \passthrough{\lstinline!public String toString()!}.
  \item
    Il default di \passthrough{\lstinline!Object.toString()!} stampa
    \passthrough{\lstinline!NomeClasse@IndirizzoHashcode!} (es.
    \passthrough{\lstinline!Date@5acf9800!}).
  \item
    La JVM invoca \textbf{implicitamente}
    \passthrough{\lstinline!toString()!} ogni volta che un oggetto viene
    concatenato a una stringa
    (\passthrough{\lstinline!"Oggi: " + today!}).
  \item
    Sovrascrivendo (\emph{overriding})
    \passthrough{\lstinline!public String toString()!} nella classe,
    possiamo personalizzare la stampa dell'oggetto.
  \end{itemize}
\end{itemize}

\begin{center}\rule{0.5\linewidth}{0.5pt}\end{center}

\subsubsection{2. Codice: Analisi del Codice Ufficiale del
Docente}\label{codice-analisi-del-codice-ufficiale-del-docente}

\paragraph{\texorpdfstring{1.
\href{file:///home/wally/Documenti/GitHub/Programmazione-II/25-26/Teoria-e-Esercitazioni/Codice-20260902/Lezioni05-06/SimpleDate/src/Date.java}{\texttt{SimpleDate/src/Date.java}}}{1. SimpleDate/src/Date.java}}\label{simpledatesrcdate.java}

\begin{lstlisting}[language=Java]
public class Date {
    // 1. Stato dell'oggetto: campi di istanza (visibilità default di package per ora)
    int day;
    int month;
    int year;

    // 2. Costante di classe condivisa da tutte le istanze
    static final String FORMAT = "dd/mm/yyyy";

    // 3. Costruttore: usa 'this' per distinguere campi da parametri
    public Date(int day, int month, int year) {
        this.day = day;
        this.month = month;
        this.year = year;
        verify(); // Chiamata al metodo di validazione interna
    }

    // 4. Metodo di validazione dello stato
    void verify() {
        if (year < 0 || month < 1 || month > 12)
            System.out.println("Illegal date!");
        else
            if (day < 1 || day > daysPerMonth(month))
                System.out.println("Illegal date!");
    }

    String print() {
        return day + "/" + month + "/" + year + " [" + FORMAT + "]";
    }

    // 5. Override di toString(): delega al metodo print()
    public String toString() {
        return print();
    }
    
    // 6. Metodo statico: calcolo puro senza bisogno di istanziare Date
    static int daysPerMonth(int month) {
        int days;
        switch(month) {
            case 4:
            case 6:
            case 9:
            case 11:
                days = 30; // Fall-through controllato per i mesi da 30 giorni
                break;
            case 2:
                days = 28; // Febbraio fisso a 28 (il bisestile arriverà a Lezione 07!)
                break;
            default:
                days = 31;
                break;
        }
        return days;
    }
}
\end{lstlisting}

\paragraph{Analisi: Dettagli Tecnici e Scelte del
Docente:}\label{analisi-dettagli-tecnici-e-scelte-del-docente}

\begin{enumerate}
\def\labelenumi{\arabic{enumi}.}
\tightlist
\item
  \textbf{Perché \passthrough{\lstinline!daysPerMonth(month)!} è
  \passthrough{\lstinline!static!}?} Non ha bisogno di leggere
  \passthrough{\lstinline!this.day!} o
  \passthrough{\lstinline!this.year!}: riceve un mese intero e
  restituisce i giorni. Essendo una funzione pura di utilità, appartiene
  alla classe e può essere chiamata anche prima di creare oggetti (come
  fa il \passthrough{\lstinline!main!}).
\item
  \textbf{Lo \passthrough{\lstinline!switch!} con fall-through}: I
  \passthrough{\lstinline!case 4: case 6: case 9: case 11:!} sono privi
  di istruzione \passthrough{\lstinline!break!} intermedia; l'esecuzione
  scorre intenzionalmente fino a
  \passthrough{\lstinline!days = 30; break;!}.
\item
  \textbf{Validazione con \passthrough{\lstinline!verify()!}}: In questa
  prima versione didattica, se la data è errata viene stampato a video
  \passthrough{\lstinline'"Illegal date!"'}. Notare che l'oggetto non
  valido viene comunque creato in memoria! Solo più avanti (in Lezione
  17 con le eccezioni) il docente mostrerà come impedire l'istanziazione
  lanciando un'eccezione.
\end{enumerate}

\begin{center}\rule{0.5\linewidth}{0.5pt}\end{center}

\paragraph{\texorpdfstring{2.
\href{file:///home/wally/Documenti/GitHub/Programmazione-II/25-26/Teoria-e-Esercitazioni/Codice-20260902/Lezioni05-06/SimpleDate/src/MainDate.java}{\texttt{SimpleDate/src/MainDate.java}}}{2. SimpleDate/src/MainDate.java}}\label{simpledatesrcmaindate.java}

\begin{lstlisting}[language=Java]
import java.util.Scanner;

public class MainDate {
    public static void main(String[] args) {
        // Chiamata a metodo statico direttamente da classe
        System.out.println(Date.daysPerMonth(4));

        // Allocazione istanze nello Heap
        Date today = new Date(14, 10, 2025);
        Date tomorrow = new Date(15, 10, 2025);

        System.out.println(today.print());
        System.out.println(tomorrow.print());
        System.out.println(today.toString());

        String str = new String("ciao");
        System.out.println(str.toString());

        // Acquisizione input utente con java.util.Scanner
        Scanner sc = new Scanner(System.in);
        System.out.println("Enter the day: ");
        int day = sc.nextInt();
        System.out.println("Enter the month: ");
        int month = sc.nextInt();
        System.out.println("Enter the year: ");
        int year = sc.nextInt();
        sc.close(); // Chiusura dello stream di input

        Date date = new Date(day, month, year);
        System.out.println(date.toString());
    }
}
\end{lstlisting}

\begin{center}\rule{0.5\linewidth}{0.5pt}\end{center}

\paragraph{\texorpdfstring{3.
\href{file:///home/wally/Documenti/GitHub/Programmazione-II/25-26/Teoria-e-Esercitazioni/Codice-20260902/Lezioni05-06/StringPlayground.java}{\texttt{StringPlayground.java}}}{3. StringPlayground.java}}\label{stringplayground.java}

Questo file illustra 3 concetti cardine sulle stringhe e
l'ottimizzazione della memoria:

\begin{lstlisting}[language=Java]
public class StringPlayground {
    public static void main(String[] args) {
        String empty = new String();
        String str = "str";
        System.out.println("len(empty): " + empty.length());
        System.out.println("len(\"str\"): " + str.length());
        System.out.println(str.charAt(1)); // Restituisce 't'

        // str.charAt(1) = 'c'; // COMPILE-TIME ERROR: Le String in Java sono IMMUTABILI!

        String same = new String("str");
        System.out.println("str: " + str + " same: " + same);

        // Confronto di puntatori (identità di memoria) vs contenuto semantico
        System.out.println("str ?= same: " + (str == same));           // false! Oggetti diversi nello Heap
        String verySame = str;
        System.out.println("str ?= verySame: " + (str == verySame));   // true! Puntano alla stessa cella
        System.out.println("str ?= same: " + str.equals(same));        // true! Stessi caratteri

        // StringBuilder e Fluent Interface (Chaining)
        int n = 5;
        int sum = n * (n + 1) / 2;
        StringBuilder sb = new StringBuilder();
        sb.append("sum(").append(sum).append(")["); // Notazione a catena (fluent)
        for (int i = 0; i < n; i++)
            sb.append(" ").append(i);
        sb.append(" ]");
        System.out.println(sb.toString()); // sum(15)[ 0 1 2 3 4 ]
    }
}
\end{lstlisting}

\paragraph{\texorpdfstring{Analisi: Le Finezze di
\texttt{StringPlayground}:}{Analisi: Le Finezze di StringPlayground:}}\label{analisi-le-finezze-di-stringplayground}

\begin{enumerate}
\def\labelenumi{\arabic{enumi}.}
\tightlist
\item
  \textbf{Immutabilità}: in Java l'oggetto
  \passthrough{\lstinline!String!} non può essere modificato dopo la
  creazione. Non esiste un metodo \passthrough{\lstinline!setCharAt()!}.
\item
  \textbf{\passthrough{\lstinline!==!} vs
  \passthrough{\lstinline!.equals()!}}:

  \begin{itemize}
  \tightlist
  \item
    \passthrough{\lstinline!str == same!} confronta se le due variabili
    contengono lo \textbf{stesso indirizzo di memoria}. Restituisce
    \passthrough{\lstinline!false!} perché
    \passthrough{\lstinline!same!} è stato creato con
    \passthrough{\lstinline!new String(...)!} forzando una nuova
    allocazione nello Heap.
  \item
    \passthrough{\lstinline!str.equals(same)!} confronta i caratteri uno
    a uno e restituisce \passthrough{\lstinline!true!}.
  \item
    \textbf{Regola d'esame assoluta}: per confrontare gli oggetti e le
    stringhe si usa \textbf{sempre \passthrough{\lstinline!.equals()!}},
    mai \passthrough{\lstinline!==!}!
  \end{itemize}
\item
  \textbf{\passthrough{\lstinline!StringBuilder!} e Fluent Notation}:

  \begin{itemize}
  \tightlist
  \item
    Ogni concatenazione con \passthrough{\lstinline!+!} su stringhe
    ordinarie crea un nuovo oggetto \passthrough{\lstinline!String!}
    temporaneo. Nei cicli questo genera un enorme sovraccarico di
    memoria per il Garbage Collector.
  \item
    \passthrough{\lstinline!StringBuilder!} alloca un buffer mutabile
    ridimensionabile internamente.
  \item
    Il metodo \passthrough{\lstinline!.append()!} restituisce un
    riferimento a \passthrough{\lstinline!this!} (l'istanza corrente del
    builder stesso), permettendo di concatenare le chiamate in cascata:
    \passthrough{\lstinline!sb.append(...).append(...)!} (\textbf{Fluent
    API / Method Chaining}).
  \end{itemize}
\end{enumerate}

\begin{center}\rule{0.5\linewidth}{0.5pt}\end{center}

\subsubsection{\texorpdfstring{3. Ciclo: Confronto con il tuo modulo
\texttt{es01}}{3. Ciclo: Confronto con il tuo modulo es01}}\label{ciclo-confronto-con-il-tuo-modulo-es01}

Nel tuo file
\href{file:///home/wally/Documenti/GitHub/Programmazione-II/25-26/esercizi-wally-25-26/es01/src/main/java/es01/Date.java}{\passthrough{\lstinline!es01/src/main/java/es01/Date.java!}}:
1. \textbf{Switch Moderno vs Tradizionale}: Avevi già applicato la
sintassi Java 14+ con la \emph{Switch Expression}:
\passthrough{\lstinline!java    static int daysPerMonth(int month) \{        return switch (month) \{            case 4, 6, 9, 11 -> 30;            case 2 -> 28;            default -> 31;        \};    \}!}
Questa forma è sintatticamente pulitissima e idiomatica in Java moderno.
2. \textbf{Allineamento con la Lezione 05-06}: Per essere al 100\%
fedeli alla versione del docente: - Aggiungere la chiamata a
\passthrough{\lstinline!verify()!} nel costruttore per il controllo di
validità. - Correggere il refuso nella costante di formato:
\passthrough{\lstinline!"dd/mm/yyyy"!} invece di
\passthrough{\lstinline!"dd/hh/yyyy"!}.

\begin{center}\rule{0.5\linewidth}{0.5pt}\end{center}

\begin{studynote}[Nota di Studio] Con le \textbf{Lezioni 05-06} abbiamo stabilito
l'architettura iniziale di \passthrough{\lstinline!Date!}, lo scope, i
metodi statici, le stringhe e \passthrough{\lstinline!StringBuilder!}.

Il prossimo blocco è la \textbf{Lezione 07 / Slide T07: \emph{Libraries
and String Class}}: - Evoluzione di
\passthrough{\lstinline!SimpleDate!}: introduzione dell'algoritmo
completo del \textbf{calendario Gregoriano per gli anni bisestili}
(\emph{leap year}). - Analisi della classe statica ausiliaria di calcolo
matematico \passthrough{\lstinline!java.lang.Math!}. - Metodi avanzati
della classe \passthrough{\lstinline!String!}
(\passthrough{\lstinline!substring!}, \passthrough{\lstinline!indexOf!},
\passthrough{\lstinline!split!}, \passthrough{\lstinline!trim!},
\passthrough{\lstinline!replace!}). - Wrapper classes per i primitivi
(\passthrough{\lstinline!Integer!}, \passthrough{\lstinline!Double!},
\passthrough{\lstinline!Character!}), \emph{Autoboxing} e
\emph{Unboxing}.
\end{studynote}

\textbf{Dammi conferma per procedere alla Lezione 07 / T07!}


\end{document}
